\section{Introduction}
Let $K$ be a field of characteristic 0. Locally nilpotent derivations $\delta$ of polynomial algebras $K[x_1,\ldots,x_d]$ and their kernels $\ker(\delta)$ are subjects of active investigation. Traditionally, the kernel of a derivation $\delta$ of $K[x_1,\ldots,x_d]$ is called the algebra of constants of $\delta$ and is denoted by $K[x_1,\ldots,x_d]^{\delta}$. The algebras of constants of locally nilpotent derivations play an essential role in the study of the automorphism group of $K[x_1,\ldots,x_d]$,
including the generation of $\operatorname{Aut}(K[x,y])$ by tame automorphisms, the Jacobian conjecture, in invariant theory, fourteenth Hilbert's problem and other important topics. See the books by Nowicki~\cite{No1}, van den Essen~\cite{E}, and Freudenburg~\cite{Fr2} for details. In particular, using locally nilpotent derivations,
Rentschler~\cite{R} gave an easy proof of the theorem of Jung--van der Kulk~\cite{J, K} that all automorphisms of $K[x,y]$ are tame. Another natural proof based on locally nilpotent derivations was given by Makar-Limanov~\cite{ML2}, see also the book~\cite{D1}. The most natural way to define the Nagata automorphism~\cite{Na}
\[
(x,y,z)\to\big(x-2\big(xz+y^2\big)y-\big(xz+y^2\big)^2z,y+\big(xz+y^2\big)z,z\big)
\]
is also in terms of locally nilpotent derivations, see Bass~\cite{B} and Smith~\cite{Sm}. The famous Jacobian conjecture is equivalent to several conjectures stated in the language of locally nilpotent derivations, see~\cite{E}. Several nice counterexamples to fourteenth Hilbert's problem are obtained as algebras of constants of locally nilpotent derivations, see the survey and the book by Freudenburg~\cite{Fr1, Fr2} and the survey by Nowicki~\cite{No2}. On the other hand, the well known theorem of Weitzenb\"ock~\cite{W} states that if $\delta$ is a nilpotent linear operator acting on the $d$-dimensional vector space $Kx_1\oplus \cdots\oplus Kx_d$, then the algebra of constants of the locally nilpotent derivation of $K[x_1,\ldots,x_d]$ which extends $\delta$ is a finitely generated algebra. A modern proof of the theorem is given by Seshadri~\cite{Ses}, with further simplification by Tyc~\cite{T}, see also~\cite{No1}. Clearly, the algebra of constants $K[x_1,\ldots,x_d]^{\delta}$ coincides with the algebra of invariants of the linear operator
